% explainer-source-hash: b65b60844a1e7b4a31378d5babc61291bbf8e105f40563de0308ddc7fe8ea1c6
% explainer-source-commit: b928cee68d8afd7dd077c24be1c4b2e0256110a9
% explainer-generated: 2026-07-16
% Regenerate with /explain-all (see WIRING.md). Do not edit this stamp by hand:
% it is how the toolchain knows whether this document still matches the code.
\documentclass[11pt,a4paper]{article}
\input{preamble}

\title{\textbf{PDF to Audiobook}\\[0.3em]\large A Deep Dive for the Complete Beginner}
\author{Portfolio explainer}
\date{}

\begin{document}
\maketitle
\thispagestyle{empty}

\begin{keyidea}{The whole project in one paragraph}
This is a desktop program that reads the words out of a PDF file and turns them into
an mp3 you can listen to. You give it one PDF or a whole batch, optionally check what
text it found and narrow it to a page range, pick a language and a regional accent,
and press one button. It shows real progress while it works, because it converts the
document one page at a time rather than in a single opaque request. It is about 578
lines of Python across two files.
\end{keyidea}

\tableofcontents
\newpage

\section{What this project is, and what it is not}

\subsection{The problem it solves}

Reading a long PDF on a screen is tiring. If the words could be spoken aloud, you
could listen while walking or cooking. Commercial apps do this behind a subscription.
This project does it locally with a small Python program.

\subsection{What it deliberately does not do}

Being precise about the boundaries matters more than a feature list, so here they are
up front:

\begin{itemize}
  \item It does \textbf{not} read scanned PDFs. A PDF that is really a photograph of a
        page contains no text to extract. This project does no OCR (see the jargon box
        below).
  \item It does \textbf{not} work offline. Speech synthesis is done by Google's servers.
  \item It does \textbf{not} have a volume control, because the speech library it uses
        does not offer one. The project states this plainly rather than shipping a
        slider that does nothing.
\end{itemize}

\begin{jargon}{OCR (Optical Character Recognition)}
Software that looks at a picture of text and works out which letters it shows. Turning
a scan of a page into real, selectable text needs OCR. This project has none, so a
scanned PDF yields nothing to read aloud.
\end{jargon}

\begin{jargon}{PDF (Portable Document Format)}
A file format for documents that look the same everywhere. Importantly, a PDF is not
simply text. It is a set of drawing instructions. A text-based PDF stores the actual
characters plus where to draw them, so a program can pull the characters back out. An
image-based (scanned) PDF stores only a picture, so there is nothing to pull out.
\end{jargon}

\begin{jargon}{TTS (Text to Speech)}
Software that converts written text into spoken audio. This project uses an online TTS
service, so it sends text over the internet and receives audio back.
\end{jargon}

\begin{jargon}{mp3}
A compressed audio file format. Its key property for this project: an mp3 is a long
chain of independent \emph{frames}, each carrying its own small header describing how to
decode it. This matters a great deal in Section~\ref{sec:concat}.
\end{jargon}

\section{Architecture: two files, one clean split}

The project is deliberately split into two modules. The rule is simple: all the real
work lives in \file{converter.py}, and it knows nothing about windows or buttons.
\file{main.py} is only a user interface that calls into it.

\begin{figure}[H]
\centering
\resizebox{\textwidth}{!}{
\begin{tikzpicture}[node distance=10mm]
  \node[box, minimum width=34mm] (gui) {\textbf{main.py}\\Tkinter GUI\\queue, buttons, progress};
  \node[box, minimum width=34mm, below=18mm of gui] (conv) {\textbf{converter.py}\\extraction, synthesis,\\playback, CLI};
  \node[extbox, minimum width=26mm, left=16mm of conv] (cli) {Command line\\\code{python3 converter.py}};
  \node[extbox, minimum width=24mm, right=42mm of gui] (dnd) {tkinterdnd2\\(optional)};
  \node[extbox, minimum width=22mm, below=18mm of conv] (gtts) {gTTS\\speech library};
  \node[extbox, minimum width=22mm, left=12mm of gtts] (pypdf) {PyPDF2\\reads PDFs};
  \node[extbox, minimum width=24mm, right=12mm of gtts] (player) {ffplay / mpv /\\afplay / aplay};
  \node[store, minimum width=18mm, below=14mm of gtts] (google) {Google\\TTS endpoint};
  \node[store, left=14mm of pypdf] (pdf) {PDF file};
  \node[store, below=14mm of player] (mp3) {mp3 file};

  \draw[flow] (gui) -- node[lbl, right]{function calls only,\\no GUI state passed} (conv);
  \draw[dflow] (dnd) -- node[lbl, above]{drag and drop,\\degrades if absent} (gui);
  \draw[flow] (cli) -- (conv);
  \draw[flow] (conv) -- (pypdf);
  \draw[flow] (conv) -- (gtts);
  \draw[flow] (conv) -- (player);
  \draw[flow] (pypdf) -- (pdf);
  \draw[flow] (gtts) -- node[lbl, above]{HTTPS} (google);
  \draw[flow] (player) -- (mp3);
\end{tikzpicture}}
\caption{The module split. Solid arrows are hard dependencies; the dashed arrow is an
optional one that the code survives without. Note that no arrow ever runs from
\file{converter.py} up into \file{main.py}: the dependency points one way only.}
\end{figure}

\begin{keyidea}{Why the one-way split is the important design decision}
Because \file{converter.py} never imports Tkinter and never touches a widget, it can be
run with no screen at all. That is what makes the command line entry point possible, and
it is what let this document's author verify the conversion logic by running it directly
(Section~\ref{sec:verify}). A version that mixed the two could only be tested by hand,
by clicking. The README records that an earlier pass of this project was exactly that
single mixed file, and that splitting it was a deliberate fix.
\end{keyidea}

\subsection{Directory structure}

\begin{figure}[H]
\centering
\begin{forest} dirtree
[pdf-to-audiobook
  [converter.py {\rmfamily\itshape all logic, plus the CLI (218 lines)}]
  [main.py {\rmfamily\itshape Tkinter GUI only (360 lines)}]
  [requirements.txt {\rmfamily\itshape three pinned dependencies}]
  [.env.example {\rmfamily\itshape deliberately empty of variables}]
  [.gitignore]
  [LICENSE {\rmfamily\itshape PolyForm Strict 1.0.0}]
  [README.md]
  [docs
    [screenshots
      [app.png]
    ]
    [explainers
      [deep-dive.tex {\rmfamily\itshape this document}]
      [brief.tex]
    ]
  ]
]
\end{forest}
\caption{Every tracked file in the repository. There is no test directory; this is a
real gap, discussed in Section~\ref{sec:honest}.}
\end{figure}

\section{The \file{converter.py} module in detail}

\subsection{The lookup tables at the top}

Three module-level data structures configure everything else.

\code{LANGUAGES} maps ten human-readable names (\code{"English"}, \code{"Japanese"}, and
so on) to the short codes the speech library expects (\code{"en"}, \code{"ja"}). The
comment in the source is honest that this is a shortlist, not the full set the library
supports.

\code{ACCENTS} maps eight names to internet domain endings such as \code{com},
\code{co.uk} and \code{co.in}. This looks odd until you know the trick behind it, which
is worth its own box.

\begin{keyidea}{Why an accent is chosen with a website address}
The speech library talks to Google Translate's speech endpoint. Google runs that service
on many regional domains, and each region's voice has that region's accent. Asking
\code{translate.google.co.uk} gives British-accented English; asking
\code{translate.google.co.in} gives Indian-accented English. So the ``accent'' setting is
literally just a choice of which country's Google domain to send the request to. The
library calls this the \code{tld}, short for \emph{top-level domain}, the last part of a
web address.
\end{keyidea}

\code{\_PLAYER\_COMMANDS} is a list, and being a list rather than a dictionary is
meaningful: it is an ordered preference. It holds four audio players with the exact
command-line arguments each needs to play a file quietly and then exit:

\begin{center}
\begin{tabular}{@{}lll@{}}
\toprule
\textbf{Order} & \textbf{Player} & \textbf{Flags used, and why} \\
\midrule
1 & \code{ffplay} & \code{-nodisp} (no window), \code{-autoexit}, \code{-loglevel quiet} \\
2 & \code{mpv}    & \code{--no-video}, \code{--really-quiet} \\
3 & \code{afplay} & none needed (this is the macOS built-in) \\
4 & \code{aplay}  & \code{-q} for quiet (this is the Linux ALSA built-in) \\
\bottomrule
\end{tabular}
\end{center}

\subsection{The three reading functions}

\code{get\_page\_count(pdf\_path)} opens the PDF and returns how many pages it has. It
is used by the GUI to fill in the Pages column.

\code{extract\_text\_from\_pdf(pdf\_path, start\_page, end\_page)} returns all the text
of a page range as one single string, joined with newlines. It is used by the preview
window.

\code{extract\_pages\_from\_pdf(pdf\_path, start\_page, end\_page)} returns a
\emph{list}, one string per page. It is used by the conversion routine.

The last two look redundant but are not: the difference between one joined string and a
list of per-page strings is precisely what makes progress reporting possible, as
Section~\ref{sec:perpage} explains.

Both share the same page-range arithmetic, and it is worth reading closely because it
contains a real defect:

\begin{lstlisting}[language=Python,caption={The range clamping, from converter.py}]
total = len(reader.pages)
first = 1 if start_page is None else max(1, start_page)
last  = total if end_page is None else min(total, end_page)
pages_text = []
for i in range(first - 1, last):
    page_text = reader.pages[i].extract_text() or ""
    pages_text.append(page_text)
\end{lstlisting}

Line by line: \code{start\_page} and \code{end\_page} are \emph{1-based} (page one is
called 1, as a human would say it) and \emph{inclusive} (both ends are converted).
Passing nothing for either end means ``from the beginning'' or ``to the end''. The
\code{max(1, ...)} stops a caller asking for page zero or a negative page. The
\code{min(total, ...)} stops a caller running off the end of the document. The
\code{- 1} in \code{range(first - 1, last)} converts the human 1-based number to the
0-based index Python lists use.

\begin{jargon}{0-based versus 1-based counting}
Python numbers the items of a list starting at 0, so the first page is
\code{reader.pages[0]}. Humans number pages starting at 1. Any code that touches both
conventions must convert between them, and off-by-one mistakes here are so common they
have their own name: \emph{off-by-one errors}.
\end{jargon}

The \code{or ""} at the end of the extraction line is small but load-bearing.

\begin{keyidea}{Why \code{or ""} is not decoration}
PyPDF2 returns \code{None} (Python's ``no value at all'') for a page it cannot read text
from, which is the normal outcome for a scanned page. Writing
\code{"\textbackslash n".join([...])} over a list containing \code{None} raises a
\code{TypeError} and the program dies. The expression \code{page.extract\_text() or ""}
means ``use the extracted text, but if it is \code{None} or empty, use an empty string
instead''. One tiny idiom converts a crash into a blank page. The README lists this as a
challenge that was actually hit.
\end{keyidea}

\begin{honest}{Verified defect: an inverted or out-of-range page range fails with a misleading message}
The clamping above handles each end independently, and never checks that
\code{first <= last}. Ask for pages 3 to 1, and \code{range(2, 1)} is empty, so you get
no text. Ask for page 99 of a 3-page document and \code{first} is 99, \code{last} is 3,
and \code{range(98, 3)} is likewise empty.

Both cases were confirmed by running the real code against a generated 3-page PDF: each
returned an empty string, and each drove the command line entry point into
\code{ValueError: No extractable text found in '...' for the selected page range}. That
message is wrong about the cause. The document is full of extractable text; the
\emph{range} is the problem. A user could reasonably conclude their PDF was a scan.

The GUI is insulated from this by luck rather than by design: its Save Page Range button
applies \code{min}/\code{max} to the two spinboxes, so it cannot submit an inverted
range. The command line has no such guard. The honest fix is a check in
\code{convert\_pdf\_to\_audio} that raises a message naming the real problem.
\end{honest}

\subsection{\code{convert\_pdf\_to\_audio}: the heart of the project}
\label{sec:perpage}

This is the function everything else exists to serve. Its signature takes the PDF path,
an optional output directory, the language and accent codes, a \code{slow} flag, the
page range, and an optional \code{progress\_callback}.

\begin{jargon}{Callback}
A function you hand to another function so it can call you back while it works. Here the
caller passes \code{progress\_callback}, and the converter calls it after each page with
the numbers ``done'' and ``total''. The converter does not know or care what happens
inside it. The GUI passes one that posts a progress message; the command line passes one
that prints \code{Page 2/3}; a test can pass none at all, and the parameter defaults to
\code{None} so nothing is called.
\end{jargon}

\begin{figure}[H]
\centering
\resizebox{0.94\textwidth}{!}{
\begin{tikzpicture}[node distance=8mm]
  \node[box] (start) {\code{convert\_pdf\_to\_audio(...)}};
  \node[box, below=of start] (extract) {Extract the page range\\into a list of strings};
  \node[dec, below=of extract, text width=22mm] (any) {any page\\has text?};
  \node[box, right=22mm of any, fill=warnred!8, draw=warnred] (err) {raise \code{ValueError}};
  \node[box, below=12mm of any] (dir) {Work out the output folder\\(default: \file{tts\_pdf} beside the PDF)\\and create it if needed};
  \node[box, below=of dir] (open) {Open \file{<name>\_tts.mp3.part}\\for writing bytes};
  \node[dec, below=of open, text width=22mm] (loop) {more pages?};
  \node[dec, right=24mm of loop, text width=22mm] (blank) {page has\\text?};
  \node[box, below=12mm of blank, minimum width=40mm] (tts) {Send the page's text to gTTS,\\append the returned mp3 bytes\\to the open file};
  \node[box, below=12mm of tts, minimum width=40mm] (cb) {Call \code{progress\_callback(i, total)}};
  \node[box, below=12mm of loop] (rename) {\code{os.replace(.part, .mp3)}\\(atomic rename)};
  \node[box, below=of rename] (ret) {return the mp3 path};

  \draw[flow] (start) -- (extract);
  \draw[flow] (extract) -- (any);
  \draw[flow] (any) -- node[lbl, above]{no} (err);
  \draw[flow] (any) -- node[lbl, right]{yes} (dir);
  \draw[flow] (dir) -- (open);
  \draw[flow] (open) -- (loop);
  \draw[flow] (loop) -- node[lbl, above]{yes} (blank);
  \draw[flow] (blank) -- node[lbl, right]{yes} (tts);
  \draw[flow] (tts) -- (cb);
  \draw[flow] (blank.east) -- ++(14mm,0) |- node[lbl, pos=0.25, right]{no: skip\\synthesis, but\\still count it} (cb.east);
  \draw[flow] (cb.west) -- ++(-56mm,0) |- (loop.west);
  \draw[flow] (loop) -- node[lbl, right]{no} (rename);
  \draw[flow] (rename) -- (ret);
\end{tikzpicture}}
\caption{The control flow of \code{convert\_pdf\_to\_audio}. The loop body is the whole
point: one network request per page, and a progress call after every page including the
blank ones.}
\end{figure}

Four decisions in this function deserve individual attention.

\subsubsection{Fail early on an empty document}

Before creating any folder or file, the function filters the extracted pages down to
those with actual content (\code{[p for p in pages if p.strip()]}) and raises a
\code{ValueError} if none remain. \code{strip()} removes surrounding whitespace, so a
page holding only spaces and newlines correctly counts as empty. This ordering means a
document with nothing to say never leaves an empty mp3 or a stray folder behind.

\subsubsection{Where the output goes}

If the caller gives no \code{output\_dir}, the function builds one: a folder named
\file{tts\_pdf} sitting next to the source PDF. The output file is the PDF's own name
with \code{\_tts.mp3} appended, so \file{report.pdf} becomes
\file{tts\_pdf/report\_tts.mp3}.

\begin{honest}{Same range, same name: a silent overwrite}
The output name is derived only from the PDF's filename. Convert pages 1 to 5 of
\file{report.pdf}, then convert pages 6 to 10 of the same file, and the second run
overwrites the first without warning. The page range is not part of the name. For a tool
whose headline feature is page ranges, this is a sharp edge. Including the range in the
filename when one is given would cost one line.
\end{honest}

\subsubsection{The \code{.part} file and the atomic rename}

The function does not write to \file{report\_tts.mp3} directly. It writes to
\file{report\_tts.mp3.part}, and only when the whole document has been synthesized does
it call \code{os.replace} to rename that into place. A \code{try/finally} block deletes
the \code{.part} file if anything went wrong.

\begin{keyidea}{Why write to a temporary name first}
If the network drops halfway through page 7 of 20, a program writing directly to
\file{report\_tts.mp3} leaves behind a real-looking mp3 containing six and a half pages.
Nothing marks it as broken. By writing to \file{.part} and renaming only on success, the
final name either does not exist or is complete: there is no in-between state. On the
same filesystem \code{os.replace} is \emph{atomic}, meaning the operating system
guarantees it either happens completely or not at all. This is a well-known professional
pattern, and finding it in a small hobby-scale utility is a genuinely good sign.
\end{keyidea}

\subsubsection{One request per page, not one per document}

This is the design insight the whole project turns on.

\begin{keyidea}{Why per-page synthesis is what buys real progress}
The obvious implementation is \code{gTTS(text=whole\_document).save(path)}: one call, one
network request. But that call is a black box. It takes a minute, says nothing, and
finishes. There is no hook inside it to report progress, so any progress bar drawn around
it would be a lie, an animation on a timer.

Splitting the work per page changes the arithmetic. Twenty pages become twenty requests,
and the code regains control between each one. That is where \code{progress\_callback}
is called, and it is why the queue can honestly say ``Converting page 7/20''. The
progress is real because the boundaries are real.

The cost is honest too: twenty small requests are slower than one large one, and twenty
chances to hit a rate limit. The project accepts that trade for truthful progress.
\end{keyidea}

Note the loop calls \code{progress\_callback(i, total)} for \emph{every} page, but only
calls gTTS for pages that have text. A blank page advances the counter instantly. That
is the right choice: progress should track pages consumed, not requests made.

\begin{honest}{Progress counts positions in the range, not page numbers}
\code{total} is \code{len(pages)}, the number of pages \emph{in the selected range}, and
\code{i} counts from 1 within that range. So converting pages 5 to 7 of a book displays
``Converting page 1/3'', not ``page 5/7''. The fraction is correct; the labels are
misleading whenever a range starts anywhere but page 1. The status text says ``page'',
which is exactly the word that makes a reader expect the document's page number.
\end{honest}

\subsection{Joining the fragments: mp3 concatenation}
\label{sec:concat}

Each page produces its own small mp3. They must become one file. The code does the
simplest possible thing: it keeps a single file handle open and asks gTTS to write each
page's bytes into it, one after another (\code{tts.write\_to\_fp(out)}).

\begin{figure}[H]
\centering
\resizebox{\textwidth}{!}{
\begin{tikzpicture}[node distance=4mm]
  \node[box, minimum width=15mm, fill=accent!12] (a1) {frame};
  \node[box, minimum width=15mm, right=1mm of a1, fill=accent!12] (a2) {frame};
  \node[box, minimum width=15mm, right=1mm of a2, fill=accent!12] (a3) {frame};
  \node[box, minimum width=15mm, right=6mm of a3, fill=okgreen!12] (b1) {frame};
  \node[box, minimum width=15mm, right=1mm of b1, fill=okgreen!12] (b2) {frame};
  \node[box, minimum width=15mm, right=6mm of b2, fill=warnred!10] (c1) {frame};
  \node[box, minimum width=15mm, right=1mm of c1, fill=warnred!10] (c2) {frame};
  \node[box, minimum width=15mm, right=1mm of c2, fill=warnred!10] (c3) {frame};

  \draw[decorate, decoration={brace, amplitude=4pt}, draw=accent]
    ([yshift=3mm]a1.north west) -- ([yshift=3mm]a3.north east) node[midway, above=3pt, lbl]{page 1 mp3};
  \draw[decorate, decoration={brace, amplitude=4pt}, draw=okgreen]
    ([yshift=3mm]b1.north west) -- ([yshift=3mm]b2.north east) node[midway, above=3pt, lbl]{page 2 mp3};
  \draw[decorate, decoration={brace, amplitude=4pt}, draw=warnred]
    ([yshift=3mm]c1.north west) -- ([yshift=3mm]c3.north east) node[midway, above=3pt, lbl]{page 3 mp3};
  \draw[decorate, decoration={brace, mirror, amplitude=5pt}, draw=inkblue]
    ([yshift=-3mm]a1.south west) -- ([yshift=-3mm]c3.south east)
    node[midway, below=5pt, lbl]{one output file: the player simply decodes frames left to right\\and never notices the joins};
\end{tikzpicture}}
\caption{Why raw byte concatenation produces a playable mp3. Every frame carries its own
header, so a decoder needs no global index and no file-level structure.}
\end{figure}

\begin{keyidea}{Why this works, and why it is still a little informal}
Most file formats would be destroyed by this. Gluing two PNG files together does not
produce a taller image; it produces a corrupt file, because a PNG has one header
describing the whole picture. An mp3 is different: it is a stream of self-describing
frames. A decoder finds a frame header, decodes that frame, and looks for the next one.
So concatenated mp3s just play, one after another.

The caveat, which the README states itself: this relies on a property of the format
rather than on a documented promise from any specification. Doing it ``properly'' would
mean adding a dependency such as \code{pydub} or shelling out to \code{ffmpeg} to
perform the same byte copy with more ceremony. The project made a reasoned trade and
wrote down why.
\end{keyidea}

This document did not take that on trust. Section~\ref{sec:verify} reports the result of
decoding a real concatenated file.

\subsection{Finding and using an audio player}

Python's standard library has no way to play audio. So \code{find\_audio\_player} walks
\code{\_PLAYER\_COMMANDS} in order and uses \code{shutil.which} to ask ``is this program
installed on this machine?'', exactly as a shell does when you type a command. It returns
the first one found, or \code{None}.

\code{play\_audio(path)} calls it, and if nothing was found it returns \code{False}
immediately. Otherwise it runs the player with \code{subprocess.run(...)} and waits for
it to finish, then returns \code{True}.

\begin{keyidea}{A boolean return is an honesty mechanism}
\code{play\_audio} returning \code{True}/\code{False} rather than just trying and
shrugging is the whole reason the GUI can tell the user ``no player found, here is the
path, open it yourself''. The README records that an early draft simply did nothing
silently on such a machine. A function that reports its own failure lets every caller
decide what to say. The \code{check=False} in the \code{subprocess.run} call means a
player that exits with an error will not raise an exception; the project treats
launching the player as the success condition, not what the player then does.
\end{keyidea}

\subsection{The command line entry point}

Guarded by \code{if \_\_name\_\_ == "\_\_main\_\_":}, the bottom of \file{converter.py}
builds an \code{argparse} parser exposing every option: \code{-o}/\code{--output-dir},
\code{-l}/\code{--lang}, \code{-t}/\code{--tld}, \code{--slow}, \code{--start-page},
\code{--end-page} and \code{--play}. Its progress callback prints \code{Page 2/3} using a
carriage return so the line rewrites itself in place instead of scrolling.

\begin{jargon}{\code{if \_\_name\_\_ == "\_\_main\_\_"}}
A standard Python idiom meaning ``only run this when the file is executed directly, not
when another file imports it''. Without it, merely importing \code{converter} from
\file{main.py} would try to parse command line arguments and crash the GUI.
\end{jargon}

\begin{honest}{The CLI reports user mistakes as raw crashes}
Confirmed by running it: pointing the command line at a file that does not exist prints
a bare \code{FileNotFoundError} traceback, and an inverted page range prints a bare
\code{ValueError} traceback. A traceback is a developer's diagnostic, not a user
message. Wrapping the call in a \code{try/except} that prints one clear line and exits
with a non-zero status is a few lines of work and is what separates a script from a
tool. The GUI, by contrast, does catch exceptions per file and surfaces them in the
queue.
\end{honest}

\section{The \file{main.py} module: the graphical interface}

\subsection{Optional drag and drop}

The very first thing of interest is a conditional import:

\begin{lstlisting}[language=Python,caption={Optional dependency handling, from main.py}]
try:
    from tkinterdnd2 import DND_FILES, TkinterDnD
    DND_AVAILABLE = True
except ImportError:
    DND_AVAILABLE = False
\end{lstlisting}

\begin{keyidea}{Graceful degradation}
\code{tkinterdnd2} bundles a native extension and does not install cleanly everywhere. If
it is missing, this code sets a flag instead of crashing. Later, the flag decides three
things: whether the window registers a drop target, which subtitle text is shown (the
fallback version says plainly that drag and drop is unavailable and why), and whether the
main window is a \code{TkinterDnD.Tk} or a plain \code{tk.Tk}. The feature disappears;
the app does not. This is called \emph{graceful degradation}, and the honest subtitle is
a nice touch: the user is told what is missing rather than left wondering.
\end{keyidea}

Note for anyone reading this document as verification: \code{tkinterdnd2} is \emph{not}
installed on the machine where these checks were run, so the drag-and-drop path is
unexercised here. The fallback path is the one that ran.

\subsection{\code{QueueItem}: one row of state}

A tiny class holding everything known about one queued PDF: its \code{path}, its
\code{name} (just the filename, for display), its \code{page\_count}, the chosen
\code{start\_page} and \code{end\_page}, a \code{status} string, and the
\code{output\_path} once converted.

Its constructor calls \code{get\_page\_count} inside a \code{try/except} that catches
every exception and falls back to \code{page\_count = 0}, so one unreadable PDF cannot
stop a batch from being queued.

\begin{honest}{The bare \code{except Exception} hides why a PDF failed}
\code{QueueItem.\_\_init\_\_} catches every possible error and records only ``0 pages''.
An encrypted PDF, a corrupt file, and a file that is not a PDF at all are collapsed into
the same silent outcome: the queue shows \code{?} pages and range \code{n/a}. The user is
never told which of those happened, and the real exception is discarded, not logged. The
failure only surfaces later, at conversion time, as the misleading ``No extractable
text'' message from Section~\ref{sec:perpage}. Keeping the exception text and showing it
in the Status column would cost almost nothing.
\end{honest}

\begin{figure}[H]
\centering
\begin{tikzpicture}
  \node[box] (q) {\code{"Queued"}};
  \node[box, below=14mm of q] (c) {\code{"Converting..."}};
  \node[box, below=16mm of c] (p) {\code{"Converting page N/M"}};
  \node[box, below=16mm of p, fill=okgreen!12, draw=okgreen] (d) {\code{"Done"}};
  \node[box, right=34mm of c, fill=warnred!10, draw=warnred] (f) {\code{"Failed: <error>"}};

  \draw[flow] (q) -- node[lbl, right]{\code{file\_start}} (c);
  \draw[flow] (c) -- node[lbl, right]{first \code{page\_progress}} (p);
  \draw[flow] (p) to[loop left] node[lbl, left]{each page:\\\code{page\_progress}} (p);
  \draw[flow] (p) -- node[lbl, right]{\code{file\_done}\\(sets \code{output\_path})} (d);
  \draw[flow] (c) -- node[lbl, above]{\code{file\_error}} (f);
  \draw[flow] (p.east) -- node[lbl, below right, pos=0.6]{\code{file\_error}} (f.south);
\end{tikzpicture}
\caption{The status of a single \code{QueueItem} over its life. These are plain strings
written straight into the Status column, not an enumeration; nothing in the code
constrains them.}
\end{figure}

\subsection{The widget layout}

\code{App.\_build\_widgets} constructs the window top to bottom: a title, the honest
subtitle, the queue table, a row of five buttons, a Voice options panel, the Convert All
button with the overall progress bar beside it, and a status label at the bottom.

The queue is a \code{ttk.Treeview}, Tk's table widget, with a tree column for the
filename and three value columns: Pages, Range, Status.

\begin{figure}[H]
\centering
\begin{forest} dirtree
[{root window (980x720, min 880x640)}
  [Label {\rmfamily\itshape ``PDF to Audiobook''}]
  [Label {\rmfamily\itshape subtitle, text depends on DND\_AVAILABLE}]
  [queue\_frame
    [Treeview {\rmfamily\itshape File / Pages / Range / Status; double-click opens preview}]
    [Scrollbar]
  ]
  [btn\_frame
    [Add Files] [Remove Selected] [Clear Queue]
    [Preview / Page Range] [Play Selected Output]
  ]
  [options\_frame {\rmfamily\itshape LabelFrame ``Voice options''}
    [Combobox {\rmfamily\itshape Language, readonly}]
    [Combobox {\rmfamily\itshape Accent (TLD), readonly}]
    [Checkbutton {\rmfamily\itshape Slow speech}]
  ]
  [convert\_frame
    [Button {\rmfamily\itshape Convert All}]
    [Progressbar {\rmfamily\itshape overall, determinate}]
  ]
  [status\_label {\rmfamily\itshape one-line result message}]
]
\end{forest}
\caption{The widget tree. Both comboboxes are \code{state="readonly"}, so the user picks
from the list and cannot type a nonsense value.}
\end{figure}

\begin{honest}{The two voice dropdowns are independent, so nonsense combinations are reachable}
Language and Accent are separate controls with no relationship in the code:
\code{lang = LANGUAGES[...]} and \code{tld = ACCENTS[...]} are read independently in
\code{start\_conversion}. Nothing stops a user selecting Japanese with the Ireland accent
or Hindi with the South Africa accent. The regional domains carry meaningful accents for
English and very little for anything else, so most of the 80 available combinations are
meaningless. The UI presents all of them as equally valid. Greying out accents for
non-English languages, or labelling the control ``English accent'', would set honest
expectations.

Smaller and in the same area: \code{ACCENTS} lists both \code{"Default (.com)"} and
\code{"US (.com)"}, which map to the identical value \code{"com"}. Two menu entries do
exactly the same thing.
\end{honest}

\subsection{Queue management}

\code{add\_files} opens a native multi-select dialog filtered to \code{*.pdf}.
\code{\_on\_drop} handles a drop event, splits the platform-specific string of paths with
\code{self.root.tk.splitlist}, and keeps only names ending in \code{.pdf}, reporting
plainly if the drop contained none.

Both funnel into \code{\_add\_paths}, which de-duplicates against paths already queued by
building a set first, counts what it added, and reports either ``N file(s) queued'' in
green or ``Those file(s) are already in the queue'' in red.

\code{\_refresh\_tree} rebuilds the whole table from \code{self.items}. It is the single
place the display is regenerated, and it sets each row's \code{iid} (Tk's identifier for
a row) to the item's position in the list, as a string.

\begin{honest}{Using the list position as the row identifier is the root of the bug in Section~\ref{sec:bug}}
\code{\_refresh\_tree} assigns \code{iid=str(i)}, and \code{remove\_selected} and
\code{\_selected\_item} convert that identifier straight back into a list index with
\code{int(iid)}. The identifier is therefore not stable: it is a position that changes
whenever the list changes. Remove the first of three files and the third item's
identifier silently changes from \code{"2"} to \code{"1"}. Nothing else in the program
knows that happened. A stable identity, the file path or a counter that never repeats,
would remove a whole class of error.
\end{honest}

\subsection{Threads, queues, and the event loop}

This is the most technically interesting part of the file, and it needs three concepts
first.

\begin{jargon}{Event loop}
A GUI program spends its life in a loop: wait for something to happen (a click, a
keypress), run the code for it, repeat. In Tkinter this loop is started by
\code{root.mainloop()}. Crucially, while your code runs, the loop is not looping, so
nothing gets redrawn. Any slow work done directly in a button handler freezes the window.
\end{jargon}

\begin{jargon}{Thread}
A second line of execution inside the same program, running at the same time as the
first. Putting slow work on a background thread lets the event loop keep spinning, so the
window stays alive.
\end{jargon}

\begin{keyidea}{The rule that shapes this entire design: Tkinter is not thread-safe}
Widgets may only be touched from the thread that created them. A background thread that
calls \code{label.config(...)} directly may work, may garble the display, or may crash
the process, unpredictably. So the worker thread cannot update the queue itself. It needs
a way to hand facts to the main thread and let the main thread do the drawing. That way
is a \code{queue.Queue}.
\end{keyidea}

\begin{jargon}{\code{queue.Queue}}
A list designed to be safely shared between threads. One thread \code{put}s items in,
another \code{get}s them out, and the class handles all the locking internally. It is the
standard, boring, correct answer to ``how do two threads talk to each other in Python''.
\end{jargon}

The mechanism has three parts working together:

\begin{enumerate}
  \item \code{start\_conversion} (main thread) reads the dropdown values, disables the
        Convert button, sets the progress bar's maximum to the number of files, and
        starts a \code{threading.Thread} running \code{\_convert\_worker}. It is a
        \emph{daemon} thread, meaning it will not keep the program alive if the user
        closes the window.
  \item \code{\_convert\_worker} (worker thread) loops over the files, calls
        \code{convert\_pdf\_to\_audio}, and \code{put}s messages onto
        \code{self.progress\_queue}. It touches no widget.
  \item \code{\_poll\_progress\_queue} (main thread) drains the queue, updates the
        widgets, and then uses \code{self.root.after(100, self.\_poll\_progress\_queue)}
        to schedule itself to run again in 100 milliseconds. It is started once in the
        constructor and re-arms itself forever after.
\end{enumerate}

\begin{jargon}{\code{root.after(ms, func)}}
Tkinter's timer: ``run \code{func} in \code{ms} milliseconds, from the main thread''. A
function that ends by calling \code{after} on itself creates a repeating poll. Note it
schedules \emph{once}: if a run of the function ends without reaching that line, the
chain stops permanently. That fact is the whole of Section~\ref{sec:bug}.
\end{jargon}

\begin{figure}[H]
\centering
\resizebox{\textwidth}{!}{
\begin{tikzpicture}
  % lifelines
  \node[box, minimum width=30mm] (ui) at (0,0) {Main thread\\(Tk event loop)};
  \node[store, minimum width=22mm] (q) at (6,0) {\code{progress\_queue}};
  \node[box, minimum width=30mm] (w) at (12,0) {Worker thread\\\code{\_convert\_worker}};
  \node[extbox, minimum width=24mm] (c) at (17.6,0) {\code{converter.py}\\+ gTTS};
  \draw[dashed, draw=linegrey] (ui) -- (0,-11.4);
  \draw[dashed, draw=linegrey] (q) -- (6,-11.4);
  \draw[dashed, draw=linegrey] (w) -- (12,-11.4);
  \draw[dashed, draw=linegrey] (c) -- (17.6,-11.4);

  \draw[flow] (0,-1.1) -- node[lbl, above]{Convert All: disable button, start thread} (12,-1.1);
  \draw[flow] (12,-1.9) -- node[lbl, above]{\code{put(("file\_start", index))}} (6,-1.9);
  \draw[flow] (12,-2.7) -- node[lbl, above]{\code{convert\_pdf\_to\_audio(...)}} (17.6,-2.7);
  \draw[flow] (17.6,-3.5) -- node[lbl, above]{\code{progress\_callback(1, 3)}} (12,-3.5);
  \draw[flow] (12,-4.3) -- node[lbl, above]{\code{put(("page\_progress", i, 1, 3))}} (6,-4.3);
  \draw[dflow] (0,-5.1) -- node[lbl, above]{poll every 100ms: \code{get\_nowait()}} (6,-5.1);
  \draw[dflow] (6,-5.9) -- node[lbl, above]{message} (0,-5.9);
  \node[lbl, anchor=west] at (0.2,-6.5) {main thread updates \code{item.status} and redraws the Treeview};
  \draw[flow] (17.6,-7.3) -- node[lbl, above]{\code{progress\_callback(2, 3)} ... and so on per page} (12,-7.3);
  \draw[flow] (17.6,-8.1) -- node[lbl, above]{returns mp3 path} (12,-8.1);
  \draw[flow] (12,-8.9) -- node[lbl, above]{\code{put(("file\_done", index, path))}} (6,-8.9);
  \draw[flow] (12,-9.7) -- node[lbl, above]{\code{put(("all\_done", converted, total))}} (6,-9.7);
  \draw[dflow] (6,-10.5) -- node[lbl, above]{main thread re-enables Convert All, writes the summary} (0,-10.5);
\end{tikzpicture}}
\caption{The full conversion sequence. The queue is drawn as its own lifeline for a
reason: no arrow ever runs directly from the worker to the main thread. Every fact the
worker learns reaches the widgets through the queue, which is exactly what makes the
design safe.}
\end{figure}

\subsection{The message protocol}

The two threads communicate with plain tuples whose first element names the message. This
is an informal protocol, but a real one:

\begin{center}
\begin{tabular}{@{}lll@{}}
\toprule
\textbf{Message} & \textbf{Payload} & \textbf{What the main thread does} \\
\midrule
\code{file\_start}    & index                & status becomes \code{"Converting..."} \\
\code{page\_progress} & index, done, total   & status becomes \code{"Converting page N/M"} \\
\code{file\_done}     & index, output\_path  & status \code{"Done"}, store path, bar +1 \\
\code{file\_error}    & index, error text    & status \code{"Failed: ..."}, bar +1 \\
\code{all\_done}      & converted, total     & re-enable the button, write the summary \\
\bottomrule
\end{tabular}
\end{center}

The polling function drains with \code{get\_nowait()} inside a \code{while True}, and the
loop is escaped by catching \code{queue.Empty}. Using the exception as the exit condition
is idiomatic Python here, not a smell.

The worker wraps each file in its own \code{try/except Exception}, so one failed PDF
sends a \code{file\_error} and the batch continues to the next file. The overall bar
advances for failures as well as successes, which is right: the bar tracks files
\emph{attempted}. At the end, \code{all\_done} carries the counts, and the summary
distinguishes all-succeeded, some-failed, and all-failed.

One subtlety worth pointing out, because it is a classic Python trap that this code gets
right:

\begin{lstlisting}[language=Python,caption={The default-argument capture, from main.py}]
def on_page(done, total, index=index):
    self.progress_queue.put(("page_progress", index, done, total))
\end{lstlisting}

\begin{keyidea}{Why \code{index=index} is necessary and not redundant}
A function defined inside a loop that refers to the loop variable does not remember the
value the variable had when the function was created; it looks the variable up when it is
\emph{called}. By then the loop has moved on, and every callback would report the last
index. Binding it as a default argument (\code{index=index}) captures the value at
definition time. The code does this correctly. It is a small detail that a lot of
otherwise-good Python gets wrong.
\end{keyidea}

\subsection{The verified concurrency bug}
\label{sec:bug}

\begin{honest}{Verified, serious: pressing Clear Queue or Remove Selected during a conversion permanently freezes the interface}
\textbf{This was reproduced by running the real \code{App} class against a real Tk window,
not deduced from reading.}

The mechanism. \code{start\_conversion} disables only the Convert All button. Add Files,
Remove Selected and Clear Queue all stay live while the worker runs. \code{clear\_queue}
sets \code{self.items = []}. The worker, meanwhile, is iterating the \emph{old} list
object and keeps posting messages carrying indices into it. The next poll executes
\code{self.items[index]} against the now-empty list and raises \code{IndexError}.

The consequence is what makes this severe rather than cosmetic. The only \code{except} in
\code{\_poll\_progress\_queue} catches \code{queue.Empty}. An \code{IndexError} propagates
out of the function, so execution never reaches the final line,
\code{self.root.after(100, self.\_poll\_progress\_queue)}. The self-rescheduling chain
breaks and is never restarted. The polling loop is dead for the rest of the session.

Observed, with a withdrawn Tk window and a captured exception handler:
\begin{itemize}
  \item \code{errors seen: [('IndexError', 'list index out of range')]}
  \item A message posted afterwards was never consumed: the queue never drains again,
        confirming the loop is gone.
  \item \code{convert button state after crash: disabled}, and it stayed \code{disabled}
        even after an \code{all\_done} message was queued, because nothing is left alive
        to read it.
\end{itemize}

So the user sees: a frozen queue that never updates again, a stalled progress bar, and a
Convert All button greyed out forever. Only restarting the app recovers. The worker
thread, meanwhile, carries on to completion in the background and writes its mp3 files
successfully, which the user has no way of learning from the interface.

\code{remove\_selected} is the same fault with a quieter face: it removes items and
re-indexes the survivors, so in-flight messages either hit the wrong row (a status
written onto an unrelated file) or run off the end and kill the loop identically.

Two independent fixes, either sufficient: disable the queue-mutating buttons for the
duration of the conversion, exactly as Convert All already is; or make messages carry a
stable identifier (the file path, or the \code{QueueItem} itself) instead of a list
position. Wrapping the poll body in a broad \code{try/except} that still re-arms the
timer would prevent the permanent freeze, but it treats the symptom. This is the single
most valuable finding in this document.
\end{honest}

\subsection{Preview and playback}

\code{open\_preview} builds a \code{Toplevel} (a second window) with two spinboxes for the
range, a text area, a Load Preview Text button and a Save Page Range button. When the
range yields nothing it shows an explicitly helpful message naming the likely cause: a
scanned or image-only page. \code{save\_range} applies \code{min}/\code{max} to the two
spinbox values, which as noted is what keeps the GUI clear of the inverted-range defect.

\code{play\_selected} refuses politely if nothing is selected or the file is not converted
yet, then runs \code{play\_audio} on \emph{another} background thread, because that call
blocks until the audio finishes. If no player was found, it hops back to the main thread
with \code{self.root.after(0, ...)} to show the message box. That detail is the
thread-safety rule from earlier being applied correctly a second time, in a place where
it would have been easy to forget.

\section{Dependencies}

\begin{center}
\begin{tabular}{@{}llp{7.4cm}@{}}
\toprule
\textbf{Package} & \textbf{Pin} & \textbf{Role} \\
\midrule
PyPDF2      & 3.0.1  & Reads PDFs; provides \code{PdfReader}, \code{.pages} and \code{.extract\_text()}. \\
gTTS        & 2.5.4  & Sends text to Google Translate's speech endpoint and returns mp3 bytes. \\
tkinterdnd2 & 0.6.2  & Optional. Adds drag and drop. Bundles a native Tk extension, which is why it is treated as optional. \\
\bottomrule
\end{tabular}
\end{center}

Tkinter itself is not listed because it ships with Python. All three are pinned to exact
versions with \code{==}, which is the right call for an application: it means a fresh
install today behaves like the one the author tested.

\begin{keyidea}{The dependency that defines the project's limits}
gTTS is not a speech engine. It is a thin client for a Google web endpoint. Everything
that follows from that is unavoidable: the tool needs an internet connection, it is
subject to rate limits, it depends on an endpoint Google never promised to keep stable,
and it is the reason there is no volume control. The README says all of this in bold
rather than burying it. \code{.env.example} contains no variables at all, and says so:
there is no API key, because the endpoint needs none.
\end{keyidea}

\section{What was actually run to verify this document}
\label{sec:verify}

Nothing in the sections above was taken on trust where it could be executed. On the
machine where this was written, PyPDF2 3.0.1, gTTS and a working network were available;
\code{tkinterdnd2} was not. A 3-page PDF and a 2-page text-free PDF were generated for
the purpose.

\begin{center}
\begin{tabular}{@{}p{6.6cm}p{8.2cm}@{}}
\toprule
\textbf{Claim tested} & \textbf{Result} \\
\midrule
\code{get\_page\_count} & Returned 3 for the 3-page PDF. \\
Page range extraction   & Range 2 to 3 returned exactly pages 2 and 3, inclusive at both ends as documented. \\
Blank document handling & The text-free PDF raised the expected \code{ValueError} rather than crashing or writing an empty mp3. \\
Inverted range 3 to 1   & Returned an empty string, then the misleading \code{ValueError}. Defect confirmed. \\
Out-of-range start 99   & Same as above. Defect confirmed. \\
\code{find\_audio\_player} & Found \code{ffplay} first, as the ordering intends. \\
Full conversion         & Succeeded. Progress fired 1/3, 2/3, 3/3. Output 77{,}952 bytes. \\
\textbf{mp3 concatenation} & \code{file} reports \code{MPEG ADTS, layer III, v2, 64 kbps, 24 kHz, Monaural}. \code{ffmpeg} decoded the whole file to WAV with \textbf{no errors}. The concatenation claim holds. \\
\textbf{Range vs full duration} & Full 3 pages: \textbf{9.744 s}. Pages 2 to 3: \textbf{6.576 s}. That is exactly 3:2, or 3.248 s per page. The page range genuinely controls what is synthesized. \\
GUI freeze bug          & Reproduced against a real \code{App} and a real Tk root. See Section~\ref{sec:bug}. \\
CLI error handling      & Missing file and inverted range both print raw tracebacks. \\
\bottomrule
\end{tabular}
\end{center}

The duration ratio deserves a note: it is not merely ``smaller, therefore probably
right''. Three equal-length pages producing 9.744 s and two of them producing 6.576 s is
an exact 3:2 with a per-page figure that matches to the millisecond. That is strong
evidence the concatenation drops nothing and duplicates nothing.

\section{Honest findings, collected}
\label{sec:honest}

Everything below is stated in full in its own section above; this is the summary a reader
should leave with, worst first.

\begin{enumerate}
  \item \textbf{The queue-mutation freeze (Section~\ref{sec:bug}).} Verified by
        execution. Clear Queue or Remove Selected during a conversion kills the progress
        polling loop permanently and leaves Convert All disabled forever. This is a
        bricked interface from a single ordinary click, and it is the one thing here
        worth fixing first.
  \item \textbf{Row identity is a list position.} The direct cause of the above, and a
        source of silently mislabelled rows even when it does not crash.
  \item \textbf{Inverted and out-of-range page ranges produce a message blaming the
        document.} Verified. Only the GUI's \code{min}/\code{max} accident keeps users
        away from it; the CLI is fully exposed.
  \item \textbf{The CLI prints tracebacks at users.}
  \item \textbf{Silent overwrite when converting two different ranges of one PDF.} The
        output name ignores the range.
  \item \textbf{Progress labels say ``page N/M'' but count positions within the range,}
        so a range starting at page 5 reports page 1.
  \item \textbf{\code{QueueItem} swallows every exception} and reduces encrypted, corrupt
        and non-PDF files to an indistinguishable \code{?}.
  \item \textbf{Language and accent are independent}, so most of the 80 combinations the
        UI offers are meaningless, and \code{ACCENTS} contains a duplicate.
  \item \textbf{No automated tests at all.} There is no test directory. Every check in
        Section~\ref{sec:verify} was written from scratch for this document and thrown
        away. The clean module split means a pytest suite over \file{converter.py} would
        be straightforward; its absence is a choice, not an obstacle. The README lists
        this itself under ``What I'd do differently''.
  \item \textbf{Offline it does nothing.} Correctly and prominently documented, but it
        remains the largest functional limit.
  \item \textbf{Raw mp3 concatenation rests on a format property, not a specification
        guarantee.} It verified clean here. The README already flags it.
\end{enumerate}

\begin{keyidea}{The fair verdict}
The bones of this project are better than its size suggests. The one-way module split,
the atomic \code{.part} rename, the queue-based thread communication, the
default-argument capture, the boolean-returning \code{play\_audio}, and the refusal to
ship a fake volume slider are all decisions a careful engineer makes. The README is
unusually honest, including about the things this document confirms.

The weaknesses cluster in one place: the boundary between the GUI's data model and the
worker thread, where positions are used as identities and only one of four dangerous
buttons is disabled. That is where the only severe bug lives, and it is a contained,
well-understood fix rather than a design flaw. The other findings are polish.
\end{keyidea}

\section{Reading the code in the right order}

For someone opening the repository for the first time:

\begin{enumerate}
  \item \file{converter.py}, the three lookup tables, to see what is configurable.
  \item \code{extract\_pages\_from\_pdf}, the shortest function that matters.
  \item \code{convert\_pdf\_to\_audio}, top to bottom. This is the project.
  \item The \code{\_\_main\_\_} block, to see the whole thing driven with no GUI at all.
  \item \file{main.py}, \code{QueueItem} first, then \code{\_build\_widgets}.
  \item \code{start\_conversion}, \code{\_convert\_worker} and
        \code{\_poll\_progress\_queue} as one unit. They only make sense together, and
        they are where the interesting bug lives.
\end{enumerate}

\end{document}
